% Zone „Das kennst du schon“ – aus bank/flaecheninhalt-und-volumen-im-raum/zone.jsonl (zusammenbau v0.1)
\einheitenkopf[zone][Kennst du schon]{Das kennst du schon}
\setcounter{aufgabe}{0}

%% TODO Titel: Ich-kann-Satz fehlt in der Bank (Platzhalter: Kettenname) – Nr. 1
%% TODO Anweisung: ein Satz über der ersten Teilaufgabe fehlt in der Bank – Nr. 1
\begin{aufgabe}{Flächenformeln der ebenen Figuren (Dreieck, Trapez, Parallelogramm, Raute/Drachen, Kreis) kennen und anwenden}
\begin{teile}
\teil Dreieck: Grundseite 8\,cm, Höhe dazu 5\,cm. Flächeninhalt? \leerfeld[cm²]
\teil Trapez: parallele Seiten 6{,}5\,cm und 3{,}5\,cm, Höhe 4\,cm. Flächeninhalt? \leerfeld[cm²]
\end{teile}
\end{aufgabe}

%% TODO Titel: Ich-kann-Satz fehlt in der Bank (Platzhalter: Kettenname) – Nr. 2
%% TODO Anweisung: ein Satz über der ersten Teilaufgabe fehlt in der Bank – Nr. 2
\begin{aufgabe}{Volumen- und Oberflächenformeln der Elementarkörper (Prisma, Pyramide, Zylinder, Kegel) kennen}
\begin{teile}
\teil Quader: 5\,cm lang, 4\,cm breit, 3\,cm hoch. Volumen? \leerfeld[cm³]
\teil Zylinder: Radius 2\,cm, Höhe 10\,cm. Volumen? Runde auf eine Stelle. \leerfeld[cm³]
\end{teile}
\end{aufgabe}

%% TODO Titel: Ich-kann-Satz fehlt in der Bank (Platzhalter: Kettenname) – Nr. 3
%% TODO Anweisung: ein Satz über der ersten Teilaufgabe fehlt in der Bank – Nr. 3
\begin{aufgabe}{Verbindungsvektoren bilden, Beträge berechnen, Mittelpunkte bestimmen}
\begin{teile}
\teil Verbindungsvektor von $P(1 | 2 | 3)$ nach $Q(4 | 6 | 3)$? \leerfeld
\teil Mittelpunkt der Strecke von $P(-2 | 5 | 1)$ nach $Q(4 | -1 | 6)$? M( \leerfeld | \leerfeld | \leerfeld )
\end{teile}
\end{aufgabe}

%% TODO Titel: Ich-kann-Satz fehlt in der Bank (Platzhalter: Kettenname) – Nr. 4
%% TODO Anweisung: ein Satz über der ersten Teilaufgabe fehlt in der Bank – Nr. 4
\begin{aufgabe}{Skalarprodukt und „null heißt senkrecht“}
\begin{teile}
\teil Skalarprodukt $(2 | 1 | 0) \cdot (1 | -2 | 5)$? \leerfeld
\teil Stehen $(3 | -2 | 1)$ und $(2 | 4 | 2)$ senkrecht aufeinander? \leerfeld
\end{teile}
\end{aufgabe}

%% TODO Titel: Ich-kann-Satz fehlt in der Bank (Platzhalter: Kettenname) – Nr. 5
%% TODO Anweisung: ein Satz über der ersten Teilaufgabe fehlt in der Bank – Nr. 5
\begin{aufgabe}{Satz des Pythagoras für Seitenhöhen und fehlende Katheten, Strahlensätze und Ähnlichkeit für Streckfaktoren}
\begin{teile}
\teil Rechtwinkliges Dreieck, Katheten 6\,cm und 8\,cm. Hypotenuse? \leerfeld[cm]
\teil Rechtwinkliges Dreieck, Hypotenuse 2{,}5\,m, eine Kathete 1{,}5\,m. Andere Kathete? \leerfeld[m]
\end{teile}
\end{aufgabe}

%% TODO Titel: Ich-kann-Satz fehlt in der Bank (Platzhalter: Kettenname) – Nr. 6
%% TODO Anweisung: ein Satz über der ersten Teilaufgabe fehlt in der Bank – Nr. 6
\begin{aufgabe}{Gleichungen aus Formeln lösen (nach einer Größe umstellen, Wurzelgleichungen, quadratische Terme, Prozentrechnung für Anteile)}
\begin{gleichungsraster}[2]
\gl{\text{Löse $\tfrac{1}{2} \cdot 6 \cdot h = 18$ nach $h$.}} &
\gl{\text{Löse $\tfrac{1}{3} \cdot 27 \cdot h = 45$ nach $h$.}} \\
\end{gleichungsraster}
\end{aufgabe}

%% TODO Titel: Ich-kann-Satz fehlt in der Bank (Platzhalter: Kettenname) – Nr. 7
%% TODO Anweisung: ein Satz über der ersten Teilaufgabe fehlt in der Bank – Nr. 7
\begin{aufgabe}{Flächenformeln der ebenen Figuren (Dreieck, Trapez, Parallelogramm, Raute/Drachen, Kreis) kennen und anwenden – Fehler finden}
\begin{teile}
\teil Ein Dreieck hat die Grundseite 10\,cm, eine schräge Seite 5\,cm und die Höhe 4\,cm zur Grundseite. Lisa rechnet: \rechnung{A &= \tfrac{1}{2} \cdot 10 \cdot 5 \\ &= 25\,\mathrm{cm}^2} Finde den Fehler und rechne richtig.
\end{teile}
\end{aufgabe}

%% TODO Titel: Ich-kann-Satz fehlt in der Bank (Platzhalter: Kettenname) – Nr. 8
%% TODO Anweisung: ein Satz über der ersten Teilaufgabe fehlt in der Bank – Nr. 8
\begin{aufgabe}{Flächenformeln der ebenen Figuren (Dreieck, Trapez, Parallelogramm, Raute/Drachen, Kreis) kennen und anwenden – selbst rechnen}
\begin{teile}
\teil Dreieck: Grundseite 12\,cm, schräge Seite 10\,cm, Höhe zur Grundseite 8\,cm. Flächeninhalt? \leerfeld[cm²]
\end{teile}
\end{aufgabe}
